\label{sec:experiments}

\paragraph{Data.}
We use WorldTrace trajectories for the Detroit core region. To focus on trip-scale corridor choice, we operate on \emph{segment-level} trajectories (full segments) rather than short sliding windows. Each trajectory is resampled to a fixed length $F=256$ for training and evaluation. Start times are stored as Unix epoch seconds, enabling time-of-day and day-of-week conditioning.

\paragraph{Temporal integrity.}
Temporal conditioning is only meaningful if $t$ is a real timestamp. We verified that the underlying segment parquet contains valid Unix epoch times and regenerated the segment-level dataset with epoch \texttt{start\_t} (instead of window-offset indices). This enables aligning time-dependent semantics such as POI open-hours.

\paragraph{Semantic context.}
We construct raster channels on a shared Detroit grid, including road hierarchy (tier-road), POI-related features (SafeGraph), and temporal encodings derived from $t$. In this draft, we use these semantics primarily for diagnostic audits and as planned conditioning inputs to the corridor selector.

\paragraph{Baselines.}
We compare against representative families:
(i) \textbf{Shortest path / $k$-shortest paths} on the road graph (map-aware but not behavioral);
(ii) \textbf{L2 regression} that deterministically predicts a trajectory (often collapses toward a straight line);
(iii) \textbf{End-to-end diffusion} that generates continuous trajectories without an explicit corridor decision;
(iv) \textbf{Map-free CascadeTraj} with a waypoint-based decision stage and residual diffusion execution.

\paragraph{Metrics.}
We report both quantitative and qualitative evidence:
\textbf{(a) Waypoint audit} measures the distance between decision-stage waypoints and GT corridors (Table~\ref{tab:wp-audit}).
\textbf{(b) Corridor diversity} evaluates whether generated samples cover multiple GT corridors for the same OD (Fig.~\ref{fig:diagnosis-mapfree}).
\textbf{(c) Structural realism} (planned) includes on-road rate and distance-to-road statistics (e.g., \% of points within a tolerance of the road graph).
\textbf{(d) Behavioral realism} (planned) measures \emph{context-conditioned distribution matching}.
We bucket trajectories by temporal context (e.g., rush hour vs.\ midnight; weekday vs.\ weekend) and (optionally) by semantic context (e.g., high vs.\ low POI activity or road-tier exposure).
For each bucket, we compute a GT corridor distribution over $\mathcal{C}$ and compare it to the generated distribution using Jensen--Shannon divergence (JSD) or earth mover's distance (EMD).
As a counterfactual test (semantic intervention), disabling or zeroing the semantic/temporal inputs should increase this distribution gap if semantics are truly causal for corridor choice.

\paragraph{Visualization protocol.}
Following the principle that metrics can be misleading, we pair every gate with case-level visualizations: GT corridor structure, generated samples overlaid with a road mask, and intervention-style ablations (none/zeros/shuffle) for semantic channels.
